\documentclass[conference]{IEEEtran}
\IEEEoverridecommandlockouts
\usepackage{cite}
\usepackage{amsmath,amssymb,amsfonts}
\usepackage{algorithmic}
\usepackage{graphicx}
\usepackage{textcomp}
\usepackage{xcolor}
\usepackage{booktabs}
\def\BibTeX{{\rm B\kern-.05em{\sc i\kern-.025em b}\kern-.08em
    T\kern-.1667em\lower.7ex\hbox{E}\kern-.125emX}}

\begin{document}

\title{Calibrated probabilistic forecasting of football match
outcomes using machine learning and advanced feature engineering}

\author{\IEEEauthorblockN{Edelweiss Gevorgyan}
\IEEEauthorblockA{\textit{College of Science and Engineering} \\
\textit{American University of Armenia}\\
Yerevan, Armenia \\
edelveys\_gevorgyan@edu.aua.am}
\and
\IEEEauthorblockN{Gagik Khalafyan}
\IEEEauthorblockA{\textit{College of Science and Engineering} \\
\textit{American University of Armenia}\\
Yerevan, Armenia \\
gkhalafyan@aua.am}
}

\maketitle

\begin{abstract}
Football match prediction has been gaining increasing popularity as a classification problem in recent studies, yet the value of any predictive model lies in finding the right balance between correctly identifying the outcome and making sure that the predicted probabilities reliably reflect the actual likelihood of these outcomes. Rather than treating the task of predicting football match outcome purely as a classification problem, this paper discusses a machine learning framework for predicting football match results with a particular focus on both predictive accuracy and the quality of predicted probabilities, generating calibrated estimates for three outcomes: home win, draw, and away win. Historical data from La Liga matches were used to engineer a feature set capturing match context and team strength. Four models were trained and compared, including Logistic Regression, Random Forest, XGBoost, and AutoGluon, with AutoGluon achieving the highest classification accuracy of 64.3\%. Models were evaluated progressing from standard measures such as accuracy and log-loss to calibration analysis using Expected Calibration Error, Brier Score, and Ranked Probability Score, benchmarked against bookmaker-implied probabilities with an ECE of 0.036. Temperature Scaling reduced XGBoost's ECE from 0.048 to 0.033, outperforming the bookmaker benchmark. A key finding is that classification accuracy and probability calibration are independent properties, with direct implications for model selection in any application where probability estimates drive financial or strategic decision-making.
\end{abstract}

\begin{IEEEkeywords}
football match prediction, probability calibration, machine
learning, probabilistic forecasting
\end{IEEEkeywords}

\section{Introduction}
Football match result prediction has become increasingly popular
in the fields of sports analytics and machine learning over the
past decade. Most of the existing solutions address this task as
a supervised classification task: given a set of pre-match and
in-play features, a model predicts one of two or three outcomes:
a home win or not a home win, or, in the case of a three-way
outcome, a home win, draw, or an away win. This framework has
become a useful tool for benchmarking various algorithms
performance, and leading-edge models often achieve accuracies
in the range of 60 to 65 percent on major European league
data~\cite{b1}\cite{b2}\cite{b3}.

However, framing football match prediction purely as a
classification task often results in fundamental mismatch when
the intended application uses probabilistic decision-making. In
risk assessment, betting markets, or any other context where
decisions are based not only on the predicted class label, but
also magnitude of predicted probabilities, the reliability of
these probabilities becomes one of the most important metrics
of the model.

The distinction between predicting the right outcome and
predicting a reliable probability is formally captured by the
concept of calibration. A model is considered well-calibrated
if its predicted probabilities correspond to the actual
frequencies of observed outcome. Formally, among all events
assigned predicted probability p, approximately a proportion p
of those events should actually occur. Calibration is therefore
not a property of the most likely predicted class, but a
property of the entire distributional profile, and it cannot
be inferred from the model accuracy alone.

Despite the importance, probability calibration remains
underdeveloped in football prediction. Most of the existing
studies pervasively evaluate models using accuracy, F1 score,
or simulated betting profit. While these metrics measure whether
the predicted outcome is correct, they do not explicitly measure
the reliability and trustworthiness of the predicted
probabilities. Most existing machine learning models produce
probability estimates as an architectural consequence, yet this
information is rarely subjected to calibration diagnosis. As a result, the question of whether the predicted probabilities of common football match predicting models can be trustworthy and serve as a basis for strategic decision-making remains largely unaddressed.

Addressing that gap directly, this paper's pivotal contribution
consists of a systematic framework for evaluating the
probabilistic robustness of machine learning models in football
match prediction, combining calibration as a primary performance
metric alongside classification performance. The framework is
applied to Spanish La Liga match data covering ten seasons
spanning from 2015/2016 to 2024/2025, using four models across
a gradient of architectural complexity: Logistic Regression,
Random Forest, XGBoost, and AutoGluon. Bookmaker-implied
probabilities serve as a calibration benchmark, providing a
functional point of reference that reflects the market-implied
consensus of the active forecasting markets.

In addressing this gap, this paper first compares four machine
learning models in terms of both classification metrics and
calibration diagnosis, including Expected Calibration Error
(ECE), Maximum Calibration Error (MCE), Brier Score, and Ranked
Probability Score. Secondly, three post-hoc calibration methods
including Platt Scaling, Isotonic Regression, and Temperature
Scaling are applied for the models that need calibration
correction. Then, model-implied probabilities are benchmarked
against bookmaker-implied probabilities to assess operational
relevance. In the end, the paper provides evidence that model
accuracy and model reliability are independent properties, and
that the selected model should account for both dimensions
depending on the intended application.
To this end, this paper addresses the following
research questions:
\begin{itemize}
    \item \textbf{} Do standard machine learning
    models produce well-calibrated results for football
    match outcomes, and how do their produced
    probabilities compare to market-implied probabilities
    across three outcome classes?
    \item \textbf{} Can post-hoc calibration methods
    improve the probabilistic quality of poorly
    calibrated models without degrading classification
    accuracy?
\end{itemize}

\section{Related work}
A wide range of machine learning models have been applied to
football match outcome prediction across major European Leagues.
These models include Random Forest, Support Vector Machine (SVM),
XGBoost and its expansions, Neural Networks, and Long Short-Term
Memory architectures~\cite{b1}\cite{b2}\cite{b3}. These
predictions typically integrate historical match statistics, team
ratings (FIFA, Elo), player attributes, and bookmaker odds, with
state-of-the-art models achieving accuracies in the range of 63
to 65 percent on English Premier League and international match
data. By addressing class distribution, balanced sampling
techniques have been proven to extensively increase predictive
performance, with binary classifications consistently
outperforming multiclass (3-class) predictions in multiple
experimental setups~\cite{b3}. Apart from accuracy, some studies
suggest that machine learning models can generate positive betting
returns under simulation and in certain setups even outperform
bookmaker prediction performance in terms of
accuracy~\cite{b4}. With the development of advanced team and
player performance metrics such as expected goals (xG) and dynamic
team ratings, the prediction accuracies have been consistently
bolstered relative to traditional goal-based
approaches~\cite{b5}.
However, these studies share a common limitation: model selection
and evaluation are driven entirely by classification performance,
without examining whether the predicted probabilities themselves
are reliable. Rodrigues and Pinto~\cite{b1} evaluate performance
through accuracy and profit simulation only, Rahman~\cite{b2}
focuses on predictive accuracy without any probabilistic
assessment, and Choi et al.~\cite{b3} similarly rely on accuracy,
F1 score, and AUC as primary criteria. As a result, whether the
probability outputs of these models can be trusted for
decision-making remains an open question.

Despite the expansive sense of this literature, evaluation of
model performance remains primarily classification-focused. Model
performance is typically evaluated through accuracy, F-1 score,
area under the curve, simulated return on investment (ROI), and
the probabilistic reliability of outputs is rarely examined.

An emerging but expanding area of work has shifted the attention
from pure classification accuracy to the quality of probabilistic
performance in football prediction. Calibration in different
models has been shown to be dependent on the model itself: some
models benefit from post-hoc calibration, others experience
performance downgrade, and careful hyperparameter tuning to
optimize log-loss can also yield well-calibrated results without
additional calibration correction~\cite{b6}.

Wheatcroft and Sienkiewicz~\cite{b6} evaluate seven machine
learning models across 35,733 matches from 22 European leagues,
finding that Platt Scaling improved calibration only for certain
models such as Gaussian Naive Bayes and AdaBoost, and actually
reduced performance for others, establishing that post-hoc
calibration should not be automatically applied but must be tested
on a per-model basis. Beyond calibration correction, the study
introduces a conceptually important shift: rather than optimizing
hyperparameters for log-loss, models are tuned directly for
betting profit using value betting and Kelly strategies. While
forecast skill and profit are related, the relationship is shown
to be imperfect, suggesting that model selection should account
for the intended decision context rather than a single universal
metric. Neural Networks achieved the best performance in both
log-loss and betting profitability, further supporting the
argument that architectural complexity and proper score
optimization interact in non-trivial ways.
The choice of evaluation criterion is itself consequential: the
Ranked Probability Score, widely used in sports analytics, has
been shown to be less efficient than the logarithmic score for
the identification of robust forecasting frameworks, and its
adoption of outcome proximity as a reward metric lacks theoretical
justification~\cite{b7}.

Wheatcroft~\cite{b7} compares three proper scoring rules: the
Ignorance score (logarithmic score), the Brier Score, and the
Ranked Probability Score, across two simulation experiments
involving over 39,000 football matches. The central argument is
that sensitivity to distance, the defining property of the RPS,
does not align with the fundamental purpose of scoring rules,
because once an outcome is observed, no information is revealed
about the probabilities of unrealised outcomes. The Ignorance
score, being local, consistently identifies superior forecasting
systems with fewer observations than both the Brier Score and the
RPS, particularly when differences between candidate forecasts
are subtle. From an information-theoretic perspective, the
Ignorance score is closely related to log-likelihood, inheriting
desirable inferential properties that are absent in the RPS. This
has direct implications for model comparison in sports analytics,
where near-optimal forecasting systems are common and evaluation
efficiency is critical.
Score decomposition methods show that low accuracy scores in
football prediction originate from limited resolution rather than
miscalibration, especially for draw outcomes~\cite{b11}.

Foulley~\cite{b11} extends the evaluation of probabilistic
football forecasts beyond mean scoring rules by applying Murphy's
calibration-refinement decomposition to UEFA Champions League
forecasts generated from a Poisson regression model with
ELO-based covariates, compared against bookmaker-implied
probabilities. Results show that both approaches are generally
well-calibrated, but bookmakers exhibit slightly stronger
resolution and discrimination, especially for home and away wins,
while draw forecasts consistently display the weakest resolution
across all evaluated approaches. A key finding is that most
forecast deficiencies originate not from miscalibration but from
limited resolution, meaning that models assign probabilities that
are reliable on average but insufficiently sharp to discriminate
between outcomes. This methodological contribution shifts
attention from simple accuracy comparisons toward understanding
why forecasts succeed or fail, a perspective directly relevant to
the calibration framework developed in this paper.
Bayesian approaches, in contrast, have yielded calibration
outcomes that are competitive relative to traditional
benchmarks~\cite{b8}.

Diniz et al.~\cite{b8} develop two Bayesian multinomial-Dirichlet
models for predicting match outcomes in the Brazilian first
division, operating directly on categorical outcome counts rather
than detailed covariates. The models exploit Dirichlet-multinomial
conjugacy to derive closed-form posterior predictive
probabilities, and combine home and away information through
linear opinion pooling. Empirical evaluation across 1,710 matches
shows that these models achieve calibration performance comparable
to or better than Bradley-Terry and Poisson-based benchmarks
under proper scoring rules including the Brier Score and
logarithmic score. A key contribution is the emphasis on
probabilistic evaluation rather than classification accuracy,
with the authors explicitly assessing calibration and concluding
that some competing models overestimate probabilities for frequent
outcomes. However, the framework assumes independence between
matches and relies solely on aggregate outcome counts, ignoring
dynamic team strength evolution, player-level covariates, and
contextual variables, limiting its applicability to more complex
prediction settings.

From a wider machine learning perspective, Guo et al.~\cite{b9}
show that modern neural networks tend to be overconfident even
with high classification accuracy performance. Their study shows
that Temperature Scaling, a simple one-parameter post-hoc
correction, consistently reduces the Expected Calibration Error
without harming the overall classification performance.

This overconfidence arises from increased model depth, reduced
regularization, and the use of Batch Normalization, and is
observed even when classification accuracy continues to improve,
demonstrating that accuracy and calibration can diverge
structurally rather than as an isolated empirical result.
Among all evaluated post-hoc methods including histogram binning,
isotonic regression, Platt scaling, vector scaling, and matrix
scaling, Temperature Scaling provided the most effective and
computationally efficient calibration across tasks, adjusting
softmax outputs without changing classification decisions.
This finding is directly linked to ensemble models used in
football prediction, which show similar probability performance
patterns. Synthesizing 219 studies, a large-scale systematic
review reveals that in sports analytics and especially betting
research the preference is still given to classification
accuracy, while the call for stronger probabilistic standards and
strategic decision-oriented frameworks is gaining significant
momentum~\cite{b10}.

Galekwa et al.~\cite{b10} identify recurring challenges across
this body of work, including data quality limitations, lack of
standardized evaluation metrics, overfitting, and weak
out-of-sample validation, and argue for a shift toward
decision-oriented frameworks using portfolio-based betting
strategies to better balance return and risk. Importantly, the
review highlights that many studies still prioritize
classification accuracy despite evidence that accuracy alone does
not ensure betting profitability, a finding that directly
motivates the calibration-first framework developed in this paper.

Existing calibration research in football forecasting tends to
focus on specific models, binary formulations, or
single-competition setups. Systematic comparison of post-hoc
correction methods, comprehensive multidimensional calibration
profile, and direct benchmarking against market-implied
probabilities remain rare. This paper addresses that gap directly
by treating calibration as a primary performance criterion,
applied consistently across four models on a decade of Spanish
La Liga historical data~\cite{b13}.

Despite its contributions and the depth of surveyed literature,
this study has several limitations that distinguish it from other
work. The developed framework is restricted to a single league,
Spanish La Liga, and a single season of 380 matches, which limits
the generalization of the findings across different contexts.
Furthermore, the feature set is limited to pre-match features and
does not incorporate in-play data, player-level statistics, and
the study does not employ LSTM architecture or neural networks.
These aspects are reserved for future work.

\section{Methods}
\subsection{Dataset and experimental setup}
The dataset includes historical data from ten seasons
of Spanish La Liga, spanning from season 2015/2016 to
2024/2025, yielding a dataset of 3,699 matches in total~\cite{b13}. A
strict temporal split was used for training and testing to reflect
realistic deployment conditions and prevent data leakage: all
models were trained chronologically on seasons up to and including
2023/2024, corresponding to 3,319 matches, and tested
on the held-out 2024/2025 season, including 380 matches.
The target variable was set to full-time result, represented
as one of three possible discrete outcomes: home win, draw,
and away win, with class frequencies of 45.5\%, 26\% and
28.5\% respectively in the training set, highlighting the well-
documented home advantage effect in football.

The overall implementation follows a sequential pipeline.
Raw match data is first collected and preprocessed, after which
features are constructed through domain-driven feature engineering.
Statistical selection methods, including ANOVA F-tests and VIF
analysis, are then applied to retain only significant and
independent signals. The selected feature set is used to train
four machine learning models on historical seasons, and
classification performance is evaluated on the held-out test
season using accuracy, F1 score, and log-loss. For models
identified as poorly calibrated, post-hoc calibration methods
are applied, and calibration quality is subsequently evaluated
using Expected Calibration Error, Brier Score, and Ranked
Probability Score, benchmarked against bookmaker-implied
probabilities.

\subsection{Feature engineering}
The original feature set was built through exploratory data
analysis combined with domain knowledge, leveraging theoretically
grounded indicators from analyzed literature. After
statistical selection, fifteen features across five groups were
chosen as a main feature set. Among these, bookmaker-implied
probabilities for three outcomes were included, converted from
raw odds after removing the over-round margin. The team
strength group captured the Elo rating difference between two
teams, providing a dynamic metric of team quality updated
after every match. The form group includes rolling differentials
in points earned in the last five matches, where a win
is awarded three points, a draw one point, and a loss no
points, following standard football competition rules, defensive
and attack performance, and goals scored. The tactical
performance group included expected point difference, passes
per defensive action, deep completions, and shots on target
difference. Finally, the underlying group captured longer-term
offensive and defensive tendencies, including season-average
expected goals for home and away teams. The final feature
set was formalized using two statistical methods. ANOVA
F-tests were applied to remove statistically insignificant features,
keeping only features that have statistically significant
correlation with the target variable. VIF analysis was then
conducted to eliminate multicollinear features, ensuring that
only independent signals are included in the final set. In
addition to the three-class formulation of the task, a binary
reformulation with possible outcome Home win versus Not
home win is introduced as a complementary experiment. Given
the well-documented difficulty of draw class prediction, the
three-class method was simplified to binary class to assess
whether the feature set retains strong discriminative power
when the class imbalance is reduced. The binary reformulation
is evaluated using F1-score, ROC-AUC and accuracy.

\subsection{Models}
All models were implemented in Python using the scikit-learn
library for Logistic Regression, Random Forest, and calibration
methods, the XGBoost library for gradient boosting, and the
AutoGluon framework for automated machine learning.

Four models were evaluated, representing a diverse range
of architectural complexity. As a principled baseline, Logistic
Regression with L2 regularization was used, as it directly
optimizes the log-loss score during training, which naturally
yields well-calibrated results~\cite{b6}. Random Forest was tuned
through minimum leaf size to prevent probability compression
that usually arises from vote-fraction averaging~\cite{b12}, with a
final configuration of maximum depth 10 and minimum leaf
of 20. XGBoost was tuned via systematic grid search to minimize
multi-class log-loss.
XGBoost was selected over other gradient boosting alternatives
such as LightGBM due to its well-established performance on
tabular classification tasks and its native support for
multiclass log-loss optimization, making it a natural fit for
the calibration-focused experimental setup of this paper.
AutoGluon, an automated machine learning framework, constructs
a weighted ensemble across multiple base learners automatically;
the final configuration used the best-quality preset with a
budget of 1,800 seconds.
AutoGluon was included to provide an upper bound on ensemble
performance without manual tuning, serving as a reference point
for what automated model selection and stacking can achieve
relative to individually tuned models.

\subsection{Calibration methods}
After the evaluation of models, three post-hoc calibration
techniques were applied for Random Forest and XGBoost,
each fitted using five-fold cross-validation to prevent leakage.
Platt Scaling fits a logistic regression model on raw model
scores to adjust output probabilities~\cite{b12}.
AutoGluon was excluded from external post-hoc calibration
as it natively incorporates Temperature Scaling through its
built-in calibration parameter, which automatically applies
probability calibration on the validation set during the
fitting process, making additional external correction
redundant within this experimental setup.
Isotonic Regression applies a non-parametric monotonic
transformation to predicted probabilities and is known to
require large datasets to perform reliably~\cite{b12}.
Despite this limitation, Isotonic Regression was included to
provide a comprehensive comparison of post-hoc correction
methods and to empirically confirm whether its known
unreliability on small datasets is reproduced in this
experimental setup.
Temperature Scaling, introduced by Guo et al.~\cite{b9},
divides model logits by a single learned scalar parameter $T$
before applying the softmax function:
\begin{equation}
\hat{p}_i = \frac{\exp(z_i/T)}{\sum_j \exp(z_j/T)}
\end{equation}
where $z_i$ are the raw logits for class $i$ and $T$ is the
temperature parameter optimized on a held-out validation set.
A temperature below 1 sharpens the probability distribution
and corrects systematic under-confidence, while a temperature
greater than 1 softens it, reducing overconfidence~\cite{b9}.
Temperature Scaling is considered well-suited for calibration
correction for applications where classification performance
must be preserved~\cite{b9}.

\subsection{Model evaluation}
Models were compared across two dimensions. Classification
performance was assessed based on accuracy, per-class F-1
score, and ROC-AUC for Home win/No Home win binary
reformulation of the task. For probabilistic reliability, five
calibration metrics were computed. First, Expected Calibration
Error (ECE) approximates the average difference between
the predicted probability and observed outcome frequency
across partitioned confidence intervals, where a smaller gap
indicates better calibration and zero represents a perfect match
between predicted and observed probabilities~\cite{b9}.
Maximum Calibration Error (MCE) captures the worst-case gap
across all probability bins, providing an upper bound on
miscalibration. The Brier Score measures the mean squared
difference between predicted probabilities and binary outcome
indicators, penalizing both miscalibration and poor
sharpness~\cite{b7}. The Ranked Probability Score extends this
to ordered multiclass outcomes by penalizing probability mass
assigned too far from the actual outcome~\cite{b8}. Reliability
diagrams provide a qualitative counterpart to numerical metrics
by plotting predicted probabilities against observed frequencies
graphically. Bookmaker-implied probabilities serve as the
calibration benchmark throughout all experiments, with an
observed ECE of 0.036.

\section{Results}
All models were tested on a held-out La Liga season consisting
of 380 matches. Table~\ref{tab1} summarizes the classification
results of multi-class models alongside the betting market as a
baseline. All the models consistently outperformed the bookmaker baseline, showing that the engineered feature set captures
meaningful signals beyond what is already covered in the
bookmaker odds. Among all the models, AutoGluon achieved the
highest overall accuracy and the best log-loss, followed by
Random Forest, Logistic Regression, and XGBoost in descending
order.

Logistic Regression's strong log-loss performance in comparison
with other models suggests that the nature of the feature set
is largely linear in essence. This observation is further
supported by the fact that XGBoost underperformed both Logistic
Regression and Random Forest across all classification metrics,
despite its architectural complexity.

Predicting draw was consistently the most challenging for all
the models, as shown in the F1 Draw and Draw Recall columns in
Table~\ref{tab1}. This is consistent with the wider football
analytics literature, where draw outcomes have lower base rate
and higher outcome uncertainty~\cite{b7}\cite{b9}. Home win
prediction was the strongest across all models, showcasing the
class frequency imbalance as described in the previous section.

\begin{table}[htbp]
\caption{Multi-class classification results}
\begin{center}
\begin{tabular}{|c|c|c|c|c|c|c|}
\hline
\textbf{Model} & \textbf{Acc.} & \textbf{LL} &
\textbf{F1 H} & \textbf{F1 D} & \textbf{F1 A} &
\textbf{DR} \\
\hline
AutoGluon & 0.643 & 0.831 & 0.746 & 0.451 & 0.657 & 0.469 \\
\hline
Rand. Forest & 0.637 & 0.848 & 0.753 & 0.319 & 0.640 & 0.227 \\
\hline
Log. Reg. & 0.629 & 0.845 & 0.759 & 0.256 & 0.636 & 0.175 \\
\hline
XGBoost & 0.616 & 0.857 & 0.738 & 0.197 & 0.654 & 0.134 \\
\hline
Betting BL & 0.544 & - & - & - & - & - \\
\hline
\multicolumn{7}{l}{Acc.=Accuracy, LL=Log Loss, H=Home, D=Draw,} \\
\multicolumn{7}{l}{A=Away, DR=Draw Recall, BL=Baseline}
\end{tabular}
\label{tab1}
\end{center}
\end{table}

\subsection{Binary classification}
Given the persistent difficulties to predict draw results across
all models, a binary reformulation of the task was introduced
as a complementary approach. By changing the multiclass setting
to Home Win vs. Not Home Win reformulation, the three-class
imbalance was reduced and the model can focus on the most
frequent outcome class. Under this reformulation, all four
models consistently achieved substantially higher results, as
reported in Table~\ref{tab2}. The uniformly high performance
across all models confirms that the feature set is specifically
informative for predicting home win, and suggests that the
significant portion of prediction errors are closely tied to
the hard prediction of the draw class rather than the limitation
of the feature set or model architecture.

\begin{table}[htbp]
\caption{Binary classification results (Home win vs. not home win)}
\begin{center}
\begin{tabular}{|c|c|c|c|c|c|}
\hline
\textbf{Model} & \textbf{Acc.} & \textbf{AUC} &
\textbf{F1 HW} & \textbf{F1 NH} & \textbf{vs. BM} \\
\hline
Rand. Forest & 0.753 & 0.811 & 0.78 & 0.72 & +13.8 pp \\
\hline
Log. Reg. & 0.747 & 0.817 & 0.78 & 0.71 & +13.2 pp \\
\hline
XGBoost & 0.743 & 0.808 & 0.77 & 0.70 & +12.8 pp \\
\hline
Bookmaker BL & 0.615 & - & - & - & Reference \\
\hline
\multicolumn{6}{l}{Acc.=Accuracy, HW=Home Win, NH=Not Home,} \\
\multicolumn{6}{l}{BM=Bookmaker, BL=Baseline}
\end{tabular}
\label{tab2}
\end{center}
\end{table}

\subsection{Calibration analysis}
Table~\ref{tab3} summarizes calibration results before
and after post-hoc correction for all models, benchmarked
against bookmaker-implied probabilities, with an ECE of
0.036. Table~\ref{tab4} provides a further breakdown of
per-class calibration performance before and after
Temperature Scaling.

Logistic Regression achieved the best natural calibration
of all models before any correction, with an ECE of 0.0351,
marginally better than the bookmaker benchmark. This outcome
is consistent with the known property of logistic regression
that directly optimizing log-loss during training naturally
encourages well-calibrated probability outputs~\cite{b6}\cite{b10}.
No post-hoc correction was required. Per-class analysis in
Table~\ref{tab4} confirms that draw predictions achieved
the best calibration of the three outcome classes with an
ECE of 0.0157, while away win predictions showed the largest
calibration gap before correction, a gap of 0.0549.

Random Forest showed the worst natural calibration results,
notably worse than the bookmaker, despite outperforming it
on all the classification metrics as shown in
Table~\ref{tab1}. The model had the second highest accuracy
at 0.637, yet is the worst naturally calibrated among all
models. This is confirmed by the fact that classification
accuracy and probability reliability are independent, and
can diverge substantially in certain setups. Confidence
quintile analysis confirmed that Random Forest systematically
underestimated its own probabilities in every confidence bin,
which is a known consequence of vote-fraction probability
compression in ensemble models~\cite{b12}. Among three
post-hoc methods assessed, Temperature Scaling yielded the
best results, reducing ECE from 0.0453 to 0.0362, achieving
the benchmark while also improving log-loss, Ranked
Probability Score, and Brier Score all together. Platt
Scaling, in contrast, provided only partial improvement,
and Isotonic Regression performed poorly due to the limited
size of the test data of 380 matches, consistent with its
known requirement for large calibration datasets~\cite{b12}.

As shown in Table~\ref{tab3}, XGBoost showed the largest
natural calibration error among all models, despite being
tuned to minimize log-loss during hyperparameter search.
This suggests that log-loss optimization alone does not
guarantee well-calibrated outputs for gradient boosting
models, unlike for Logistic Regression, which is naturally
well-calibrated. Confidence quintile analysis showed
systematic underconfidence across all portions, with
high-confidence predictions having the largest gaps. The
chosen Temperature Scaling parameter dropped below one,
confirming that sharpening was required rather than
softening. After calibration correction, XGBoost achieved the best
post-calibration ECE of all models, reducing ECE from 0.0480
to 0.0327, demonstrating that a model built on
bookmaker-derived features can be successfully re-calibrated
to produce reliable probability estimates, despite having
the second lowest classification accuracy as shown in
Table~\ref{tab1}. The largest
improvement for a specific class was observed for away
team win predictions. However, the improvement was not
uniform, draw calibration increased from 0.0222 to 0.0257,
suggesting that uniform temperature parameters may not
optimally correct calibration across all outcome classes.

In general, prior to any calibration correction, only
Logistic Regression achieved the benchmark level
calibration, due to its nature. All three models matched
or exceeded the market benchmark after applying Temperature
Scaling, showing that achieving well-calibrated results is
feasible with standard machine learning models. Temperature
Scaling was the most effective post-hoc calibration method,
systematically improving the probability quality, while
preserving the classification accuracy.

\begin{table}[htbp]
\caption{Calibration results before and after post-hoc correction}
\begin{center}
\begin{tabular}{|c|c|c|c|c|c|}
\hline
\textbf{Model} & \textbf{ECE B} & \textbf{ECE A} &
\textbf{Method} & \textbf{vs. BM} & \textbf{Verdict} \\
\hline
Log. Reg. & 0.0351 & 0.0351 & None & Better & Good \\
\hline
Rand. Forest & 0.0453 & 0.0362 & T.S. (0.894) & Matches & Good \\
\hline
XGBoost & 0.0480 & 0.0327 & T.S. (0.748) & Better & Good \\
\hline
Bookmaker & 0.0361 & - & - & Ref. & Ref. \\
\hline
\multicolumn{6}{l}{ECE B=Before, ECE A=After, BM=Bookmaker,} \\
\multicolumn{6}{l}{T.S.=Temperature Scaling}
\end{tabular}
\label{tab3}
\end{center}
\end{table}

\begin{table}[htbp]
\caption{Per-class ECE before and after Temperature Scaling}
\begin{center}
\begin{tabular}{|c|c|c|c|}
\hline
\textbf{Model} & \textbf{Class} &
\textbf{ECE Before} & \textbf{ECE After} \\
\hline
Log. Reg. & Home Win & 0.0347 & 0.0347 \\
\cline{2-4}
 & Draw & 0.0157 & 0.0157 \\
\cline{2-4}
 & Away Win & 0.0549 & 0.0549 \\
\hline
Rand. Forest & Home Win & 0.0558 & 0.0367 \\
\cline{2-4}
 & Draw & 0.0318 & 0.0299 \\
\cline{2-4}
 & Away Win & 0.0483 & 0.0421 \\
\hline
XGBoost & Home Win & 0.0597 & 0.0453 \\
\cline{2-4}
 & Draw & 0.0222 & 0.0257 \\
\cline{2-4}
 & Away Win & 0.0621 & 0.0272 \\
\hline
\end{tabular}
\label{tab4}
\end{center}
\end{table}

\section{Analysis and Validation}

\subsection{Key findings}
These results reveal three main findings. First, log-loss
optimization during training is a meaningful step for cal-ibration quality. As confirmed by the results for Logistic
Regression, natural minimization of log-loss requires no post-hoc correction and produces the best probability reliability,
consistent with the results of Wheatcroft and Sienkiewicz~\cite{b6},
who show that careful optimization of a proper score during
initial training can fully substitute the calibration correction
process.

Secondly, Temperature Scaling is the most recommended
post-hoc calibration method for ensemble models in this
particular experimental setup. It consistently outperformed
other post-hoc calibration methods such as Platt Scaling and
Isotonic Regression while preserving classification performance.
Its single-parameter simplicity and computational efficiency
make it well-suited for small-to-medium datasets in sports
evaluation, and Isotonic Regression, in contrast, was shown
to be unreliable~\cite{b12}.

Lastly, and most importantly, classification accuracy and
calibration quality are mutually distinct characteristics of
a machine learning model, as clearly illustrated in
Tables~\ref{tab1},~\ref{tab3}, and~\ref{tab4}. Models with
lower classification accuracy achieved the best post-hoc
calibration ECE. This dissociation confirms that model
selection must consider both dimensions separately depending
on the intended application. For strategic decision-making
applications where probability estimates drive financial
decisions, calibration quality and probability reliability
should be weighted at least as heavily as classification
accuracy~\cite{b6}\cite{b11}.

\subsection{Validation}
The results of this paper are validated through three
complementary approaches. First, a strict temporal split
was applied throughout all experiments, ensuring that no
future data was used during training and that the evaluation
reflects realistic deployment conditions. Second,
bookmaker-implied probabilities, with an observed ECE of
0.036, serve as a functional calibration benchmark,
representing the market-implied consensus of professional
forecasting markets and providing a meaningful reference
point for assessing the probabilistic quality of model
outputs. Third, five-fold cross-validation was applied
during post-hoc calibration fitting to prevent leakage
between calibration and evaluation sets. Together, these
validation approaches ensure that the reported results
are not artifacts of data leakage or overfitting, and
that the calibration improvements observed after
Temperature Scaling reflect genuine improvements in
probabilistic reliability.

\subsection{Comparison with existing work}
The classification results obtained in this paper are
consistent with the wider football analytics literature,
where leading-edge models achieve accuracies in the range
of 63 to 65 percent on major European league data~\cite{b1}\cite{b2}\cite{b3}.
AutoGluon's best accuracy of 64.3\% falls within this
established range, confirming that the engineered feature
set is competitive with state-of-the-art approaches.
The persistent difficulty of draw prediction is consistent
with the findings of Foulley~\cite{b11}, who shows that
draw outcomes consistently display the weakest resolution
and discrimination across all evaluated models and
benchmarks. The natural calibration of Logistic Regression
without any post-hoc correction is consistent with the
findings of Wheatcroft and Sienkiewicz~\cite{b6}, who
demonstrate that careful optimization of a proper score
during training can substitute calibration correction.
The failure of Isotonic Regression on the 380-match test
set directly confirms its known requirement for large
calibration datasets as established by Kull et al.~\cite{b12}.
The effectiveness of Temperature Scaling in improving
calibration without degrading classification performance
is consistent with the findings of Guo et al.~\cite{b9},
who demonstrate this property across a range of modern
machine learning models.

\subsection{Objectives and limitations}
This paper set out to address two research questions:
whether standard machine learning models produce
well-calibrated results for football match outcomes, and
whether post-hoc calibration methods can improve
probabilistic quality without degrading classification
accuracy. Both objectives were met. All four models were
systematically evaluated across classification and
calibration dimensions, and post-hoc calibration with
Temperature Scaling successfully improved the probabilistic
quality of Random Forest and XGBoost to at or beyond the
bookmaker benchmark, while preserving classification
performance in both cases.

Several limitations of this work should be acknowledged.
The framework is restricted to a single league, Spanish
La Liga, and evaluated on a single held-out season of
380 matches, which limits the generalization of the
findings across different competitions and time periods.
The feature set is limited to pre-match features and does
not incorporate in-play data or player-level statistics,
which may capture additional predictive signals. A single
global Temperature Scaling parameter was applied across
all outcome classes, yet the per-class results in
Table~\ref{tab4} show that calibration improvement was
not uniform across home win, draw, and away win outcomes,
suggesting that class-specific calibration correction
may be more appropriate. Finally, the study does not
investigate the relationship between calibration quality
and betting profitability, which remains an important
direction for future work.

\section{Conclusion}
This paper presented a machine learning framework for
calibrated forecasting of football match outcomes, with
calibration metrics as the final evaluation criterion. Four
machine learning models - Logistic Regression, Random
Forest, XGBoost, and AutoGluon - were trained on ten
seasons of Spanish La Liga data and assessed using a holistic
array of classification and calibration diagnostics,
benchmarked against betting market-implied probabilities.
The main contributions of this work are threefold: a
systematic calibration-first evaluation framework applied
consistently across four models of varying architectural
complexity, a direct benchmarking of model-implied
probabilities against professional bookmaker probabilities,
and empirical evidence that classification accuracy and
probability calibration are independent properties with
direct implications for model selection in probabilistic
forecasting applications.

On classification performance, all models outperformed the
bookmaker baseline with AutoML model AutoGluon achieving the
best results in terms of accuracy. Draw prediction remained
the central problem of outcome prediction across all models,
a finding consistent with the broader literature. The binary
reformulation of the task confirmed that the significant
prediction error is a direct consequence of the hard
predictive nature of the draw class, rather than a limitation
of model selection or architecture.

As for calibration, the results demonstrate that achieving a
well-calibrated probabilistic forecast is feasible with
machine learning models, and post-hoc correction with a
single parameter sufficiently improves calibration metrics
and achieves the professional bookmaker baseline. Logistic
Regression showed well-calibrated results without any
corrections, while XGBoost and Random Forest required
Temperature Scaling to achieve competitive performance.
Notably, XGBoost demonstrated that a model built on
bookmaker-derived features can be successfully recalibrated
to produce reliable probability estimates.
Isotonic Regression consistently worsened the model's
probability performance, which is directly connected to the
small size of the test dataset.

The most important finding of this work is that
classification accuracy and probability quality are
independent properties. A model can achieve high accuracy
while producing unreliable probabilities, and on the other
hand, a well-calibrated model can have relatively modest
accuracy. This discrepancy has direct implications: in any
context where probability estimates drive financial
decisions, such as risk assessment, betting, and
probabilistic model comparison, calibration should be
evaluated explicitly and should not be inferred from
classification performance alone.

The two research questions posed at the outset of this paper
were both addressed. Standard machine learning models were
shown to produce varying levels of calibration quality, with
only Logistic Regression achieving benchmark-level
calibration naturally, while post-hoc calibration methods
were shown to improve probabilistic quality of ensemble
models without degrading classification accuracy, directly
answering both research questions.

The primary limitations of this work include its restriction
to a single league and a single held-out season of 380
matches, the use of pre-match features only, and the
application of a single global Temperature Scaling parameter
across all outcome classes despite non-uniform per-class
calibration improvement.

Several directions for future work have been identified.
First, a single global temperature value should not be used
for every class, as the improvement may not be uniform
across all the classes. Then, this framework can be applied
to a larger leagues and seasons dataset to properly assess
the generalizability of these findings. Finally, if this
framework is used for betting strategies or bookmaker
financial decisions, the relationship between probability
quality and profitability warrants should be properly
investigated to identify positive expected returns on
investment in practice.

\section*{Acknowledgment}
The author of this paper would like to express gratitude to Gagik Khalafyan for continuous support and guidance throughout all the stages from the very beginning of this project. His expertise and insights were instrumental for developing the analytical framework and the research direction.
The author used Claude AI (Anthropic, 2025) for structural editing and grammar enhancement of the manuscript.
This work was completed as a part of the Bachelor of Science in Data Science of the College of Science and Engineering capstone requirement of the American University of Armenia.

\begin{thebibliography}{00}

\bibitem{b1} F. Rodrigues and A. Pinto, ``Prediction of football
match results with machine learning,'' 2022. [Online]. Available:
https://www.researchgate.net/publication/363472367

\bibitem{b2} M. A. Rahman, ``A deep learning framework for
football match prediction,'' \textit{Discover Applied Sciences},
Jan. 2020. [Online]. Available: https://link.springer.com

\bibitem{b3} B. S. Choi, L. K. Foo, and S.-L. Chua, ``Predicting
football match outcomes with machine learning approaches,''
\textit{MENDEL}, vol. 29, no. 2, pp. 229--240, 2023.
https://doi.org/10.13164/mendel.2023.2.229

\bibitem{b4} E. Christoffersson, ``Beating the odds: Machine
learning for football match prediction,'' 2023. [Online].
Available: https://www.diva-portal.org/smash/record.jsf?pid=diva2:1772002

\bibitem{b5} C. Herbinet, ``Using machine learning techniques to
predict the outcome of professional football matches,''
Undergraduate Project Report, Imperial College London, 2018.
[Online]. Available: https://www.imperial.ac.uk/media/imperial-college/
faculty-of-engineering/computing/public/distinguished-projects/
1718-ug-projects/Corentin-Herbinet-Using-Machine-Learning-techniques-to-predict-the-outcome-of-professional-football-matches.pdf

\bibitem{b6} E. Wheatcroft and E. Sienkiewicz, ``Calibration and
hyperparameter tuning in football forecasting with machine
learning,'' 2023. [Online]. Available:
https://www.researchgate.net/publication/352846677

\bibitem{b7} E. Wheatcroft, ``Evaluating probabilistic forecasts
of football matches: the case against the ranked probability
score,'' \textit{Journal of Quantitative Analysis in Sports},
vol. 17, no. 4, pp. 273--287, Dec. 2021.
https://doi.org/10.1515/jqas-2019-0089

\bibitem{b8} M. A. Diniz, R. Izbicki, D. Lopes, and L. E.
Salasar, ``Comparing probabilistic predictive models applied to
football,'' \textit{Journal of the Operational Research Society},
vol. 70, no. 5, pp. 770--782, May 2019.
https://doi.org/10.1080/01605682.2018.1457485

\bibitem{b9} C. Guo, G. Pleiss, Y. Sun, and K. Q. Weinberger,
``On calibration of modern neural networks,'' in \textit{Proc.
34th Int. Conf. Mach. Learn. (ICML)}, Sydney, Australia, 2017,
pp. 1321--1330. [Online]. Available:
https://arxiv.org/pdf/1706.04599

\bibitem{b10} R. M. Galekwa, J. M. Tshimula, E. G. Tajeuna,
and K. Kyandoghere, ``A systematic review of machine learning
in sports betting: Techniques, challenges, and future
directions,'' arXiv:2410.21484, Oct. 2024. [Online]. Available:
https://arxiv.org/abs/2410.21484

\bibitem{b11} J.-L. Foulley, ``More on verification of
probability forecasts for football outcomes: score
decompositions, reliability, and discrimination analyses,''
arXiv:2106.14345, Jun. 2021. [Online]. Available:
https://arxiv.org/abs/2106.14345

\bibitem{b12} M. Kull, T. Silva Filho, and P. Flach, ``Beta
calibration: A well-founded and easily implemented improvement
on logistic calibration for binary classifiers,'' in
\textit{Proc. 20th Int. Conf. Artif. Intell. Stat. (AISTATS)},
Fort Lauderdale, FL, USA, 2017. [Online]. Available:
https://proceedings.mlr.press/v54/kull17a.html

\bibitem{b13} Football-Data.co.uk, ``Football data and
results,'' [Online]. Available: https://www.football-data.co.uk

\end{thebibliography}

\end{document}